\subsection{The Structure of Modules of Finite Length}

THEOREM 2 (Krull--Remak--Schmidt). --- Let A be a ring and M an A-module of finite length.

\begin{enumerate}
\def\labelenumi{\alph{enumi})}
\item
  There exists a finite family \((\mathrm{M}_i)_{i\in \mathrm{I}}\) of indecomposable submodules of M such that \(\mathrm{M} = \bigoplus_{i\in \mathrm{I}}\mathrm{M}_i\) , and the module M is semiprimordial.
\item
  Let \((\mathrm{M}_{i})_{i\in\mathrm{I}}\) and \((\mathrm{M}_{j}^{\prime})_{j\in\mathrm{J}}\) be two finite families of indecomposable submodules of M such that \(M=\bigoplus_{i\in I}M_{i}=\bigoplus_{j\in J}M_{j}^{\prime}\) . There exist a bijection \(\sigma\) from I to J and an automorphism u of M such that we have \(u(\mathrm{M}_{i})=\mathrm{M}_{\sigma(i)}^{\prime}\) for every \(i\in I\) .
\item
  Let \(\mathrm{N}\) be a direct factor submodule of \(\mathrm{M}\) , and let \((\mathrm{M}_i)_{i\in \mathrm{I}}\) be a finite family of indecomposable submodules of \(\mathrm{M}\) with direct sum \(\mathrm{M}\) . There exists a subset \(\mathrm{J}\) of \(\mathrm{I}\) such that \(\bigoplus_{i\in \mathrm{I - J}}\mathrm{M}_i\) is supplementary to \(\mathrm{N}\) . The module \(\mathrm{N}\) is isomorphic to \(\bigoplus_{j\in \mathrm{J}}\mathrm{M}_j\) .
\item
  Let \(\mathbf{N}\) be an A-module. If there exists an integer \(d > 0\) such that the modules \(\mathbf{M}^d\) and \(\mathbf{N}^d\) are isomorphic, then the modules \(\mathbf{M}\) and \(\mathbf{N}\) are isomorphic.
\item
  Let \(\mathrm{N}\) and \(\mathrm{P}\) be A-modules of finite length. If the modules \(\mathrm{M} \oplus \mathrm{P}\) and \(\mathrm{N} \oplus \mathrm{P}\) are isomorphic, then \(\mathrm{M}\) and \(\mathrm{N}\) are isomorphic.
\end{enumerate}

A module of finite length is both Artinian and Noetherian (VIII, p.~2, Proposition 1). Moreover, for a module of finite length, being indecomposable or primordial amounts to the same (VIII, p.~31, Proposition 4). Assertion a) then follows from Proposition 3 of VIII, p.~30. Assertions b), c), and e) follow from, respectively, Corollaries 2, 5, and 3 of Theorem 1 of VIII, p.~32. Finally, assertion d) follows from Corollary 4 of VIII, p.~35 because under the assumptions of d), the module N has finite length and is therefore semiprimordial by a).

THEOREM 3. --- Let K be a commutative field, A a K-algebra, and M and N modules of finite length. Let \(K'\) be a nonzero commutative K-algebra such that the \(\mathrm{A}_{(\mathrm{K}')}\) -modules \(\mathrm{M}_{(\mathrm{K}')}\) and \(\mathrm{N}_{(\mathrm{K}')}\) are isomorphic. Then the A-modules M and N are isomorphic.

\begin{enumerate}
\def\labelenumi{\alph{enumi})}
\item
  First, suppose that the algebra \(K'\) has finite degree d over K. Then the A-module \(\mathrm{M}_{(\mathrm{K}')}\) is isomorphic to \(M^{d}\) and the A-module \(\mathrm{N}_{(\mathrm{K}')}\) is isomorphic to \(\mathbf{N}^d\) , so that the A-modules \(\mathbf{M}^d\) and \(\mathbf{N}^d\) are isomorphic. By Theorem 2, d), the A-modules \(\mathbf{M}\) and \(\mathbf{N}\) are isomorphic.
\item
  Now, suppose that the K-algebra \(K'\) is generated by finitely many elements. Choose a maximal ideal m of \(K'\) , and set \(K'' = K'/m\) . By Hilbert's Nullstellensatz (VIII, p.~462, Corollary 1 of Theorem 1), \(K''\) is a finite-degree extension of K. By extension of scalars from \(K'\) to \(K''\) , we deduce from the \(A_{(K')}\) -linear isomorphism \(M_{(K')} \to N_{(K')}\) an \(A_{(K'')}\) -linear isomorphism \(M_{(K'')} \to N_{(K'')}\) . By part a) of the proof, the A-modules M and N are isomorphic.
\item
  Finally, let us treat the general case. Let \(u: \mathrm{M}_{(\mathrm{K}')} \to \mathrm{N}_{(\mathrm{K}')}\) be an isomorphism of \(\mathrm{A}_{(\mathrm{K}')}\) -modules and \(v: \mathrm{N}_{(\mathrm{K}')} \to \mathrm{M}_{(\mathrm{K}')}\) the inverse isomorphism. Denote by \(\mathcal{E}\) the set of K-subalgebras of \(\mathrm{K}'\) generated by finitely many elements. If E is such a subalgebra, then \(\mathrm{A}_{(\mathrm{E})}\) is identified with a subring of \(\mathrm{A}_{(\mathrm{K}')}\) , and \(\mathrm{M}_{(\mathrm{E})}\) and \(\mathrm{N}_{(\mathrm{E})}\) are identified with \(\mathrm{A}_{(\mathrm{E})}\) -submodules of \(\mathrm{M}_{(\mathrm{K}')}\) and \(\mathrm{N}_{(\mathrm{K}')}\) , respectively (II, §7, No.~7, p.~306); moreover, \(\mathrm{M}_{(\mathrm{K}')}\) and \(\mathrm{N}_{(\mathrm{K}')}\) are unions of right directed families \((\mathrm{M}_{(\mathrm{E})})_{\mathrm{E} \in \mathcal{E}}\) and \((\mathrm{N}_{(\mathrm{E})})_{\mathrm{E} \in \mathcal{E}}\) , respectively. The A-modules M and N are of finite length, hence finitely generated; let S be a finite generating subset of the A-module M and T a finite generating subset of the A-module N. There exists a K-algebra \(\mathrm{E} \in \mathcal{E}\) such that we have \(u(1 \otimes s) \in \mathrm{N}_{(\mathrm{E})}\) for every \(s \in \mathrm{S}\) and \(v(1 \otimes t) \in \mathrm{M}_{(\mathrm{E})}\) for every \(t \in \mathrm{T}\) . It follows by linearity that we have \(u(\mathrm{M}_{(\mathrm{E})}) \subset \mathrm{N}_{(\mathrm{E})}\) and \(v(\mathrm{N}_{(\mathrm{E})}) \subset \mathrm{M}_{(\mathrm{E})}\) . The maps \(u\) and \(v\) then induce inverse bijections from \(\mathrm{M}_{(\mathrm{E})}\) to \(\mathrm{N}_{(\mathrm{E})}\) and from \(\mathrm{N}_{(\mathrm{E})}\) to \(\mathrm{M}_{(\mathrm{E})}\) . These bijections are clearly \(\mathrm{A}_{(\mathrm{E})}\) -linear. Thus, the \(\mathrm{A}_{(\mathrm{E})}\) -modules \(\mathrm{M}_{(\mathrm{E})}\) and \(\mathrm{N}_{(\mathrm{E})}\) are isomorphic. By part b) of the proof, the A-modules M and N are isomorphic.
\end{enumerate}

Remark. --- Let E and F be two finite-dimensional vector spaces over a commutative field K, and let \(K'\) be an extension of K. Let u be an endomorphism of E and v an endomorphism of F, and let \(u_{(K')}\) and \(v_{(K')}\) be the endomorphisms of \(E_{(K')}\) and \(F_{(K')}\) obtained by extension of scalars. It follows from Corollaries 1 and 2 of VII, §5, No.~3, p.~32 that the endomorphisms u and v are similar if and only if the endomorphisms \(u_{(K')}\) and \(v_{(K')}\) are. This also follows from Theorem 3 above applied to the algebra A = K{[}X{]} and the A-modules \(M = E_u\) and \(N = F_v\) (VII, §5, No.~1, p.~29).

